%--- E13: Section 7.2, after the paragraph on the regimes / table ------------
\begin{explained}{Burgers: why the shock forms and what happens to the interaction}
\textbf{Step 1: why inviscid Burgers forms a shock.} $u_t+uu_x=0$ is nonlinear transport: the velocity at each point is $u$ itself. Along a characteristic starting at $x=\xi$, $dx/dt=u$ and $du/dt=0$, so $u$ stays constant, $u(x(t),t)=u_0(\xi)$. For $u_0=\sin\xi$:
\[
x=\xi+t\sin\xi,\qquad u=\sin\xi .
\]
Differentiating, $\partial x/\partial\xi=1+t\cos\xi$; as long as it is positive everywhere the map $\xi\mapsto x$ is invertible and the solution is smooth. It first vanishes when $1+t\cos\xi=0$; since $\min\cos=-1$, the breaking time is $t^*=1$ at $\xi^*=\pi$. The gradient $u_x=\dfrac{\cos\xi}{1+t\cos\xi}$ has numerator $\to-1$ and denominator $\to0$ at $(t,\xi)=(1,\pi)$, so it blows up. \emph{Before $t=1$} characteristics do not cross and there is one value of $u$ per $x$. \emph{At and after} $t=1$ they compress and cross, the parametric curve folds, no single-valued classical solution exists, and the entropy solution replaces the fold with a shock (for this symmetric datum at $x=\pi$).

\medskip\textbf{Step 2: where the interaction enters.} Viscous Burgers $u_t+uu_x=\nu u_{xx}$ is split into nonlinear transport $A_{\mathrm{tr}}u=-uu_x$ and diffusion $A_{\mathrm{diff}}u=\nu u_{xx}$. With $S(t)$ the full flow and $S_1,S_2$ the single-term flows,
\[
\Eint(t)=S(t)u_0-S_1(t)u_0-S_2(t)u_0+u_0 .
\]
Before the shock all flows can be compared with smooth expansions; after it, the transport flow develops a discontinuity while the viscous solution smooths the transition over a thin layer. That difference drives everything below.

\medskip\textbf{Step 3: why the smooth expansion breaks down.} Before breaking there is a first-order expansion in $\nu$ with remainder $O(\nu^2\varepsilon^2t^3)$ (the exact coefficient depends on the definitions of the paper; the parameter $\varepsilon$ is not defined in the text and is flagged in the author's notes). A perturbative expansion is useful while the relevant solutions stay smooth. At $t^*$ the inviscid gradient is unbounded and estimates relying on bounded derivatives fail. \emph{Precision:} what blows up at $t^*$ is the perturbative expression or its estimates, not the viscous solution itself, which stays finite because positive viscosity regularizes.

\medskip\textbf{Step 4: a jump that survives $\nu\to0$.} After the shock the transport flow is discontinuous while the viscous flow has a thin smooth layer. Even if the viscous profile converges to the entropy solution, the individual pieces of the difference need not converge so that their pointwise jumps cancel; hence $\Eint$ can have a jump of size independent of $\nu$. This does not contradict smallness in $L^1$: a function of amplitude $O(1)$ supported on an interval of width $O(\nu)$ has $\|E\|_{L^\infty}=O(1)$ but $\|E\|_{L^1}=O(\nu)$. The integral measures the area under $|E|$, not the maximum height.

\medskip\textbf{Step 5: the three estimates.}
\begin{center}\small
\begin{tabular}{@{}L{4.3cm}L{3.0cm}L{6.6cm}@{}}
\toprule
Quantity & Claimed scale & Interpretation\\\midrule
$\|\Eint(t)\|_{L^1}$ at fixed $t$ & $\Theta(\nu)$ & small in total absolute magnitude at a fixed post-shock time\\
$\sup_{t\le T}\|\Eint(t)\|_{L^1}$ & $\nu|\log\nu|+O(\nu)$ & the worst case over the time interval has a logarithmic enhancement\\
$\|S-S_1S_2\|_{L^1}$ & $\Theta(\sqrt\nu)$ & a different, larger scale: the error of replacing the full evolution by a composition of separate flows\\
\bottomrule
\end{tabular}
\end{center}
$\Theta(\nu)$ means bounded above and below by positive multiples of $\nu$ as $\nu\to0$. The third quantity is a different object: it concerns $S-S_1S_2$, not the additive remainder, and must not be conflated with it. Exact constants and rates depend on the definitions of the split flows, the domain and the theorem's hypotheses; the status labels in the text say what was proved and what is conditional.

\medskip\textbf{Step 6: the broader lesson.} Three regimes: before the shock, smoothness permits perturbative expansions and derivative-based estimates; at breaking, the inviscid gradient blows up and smoothness-based estimates fail; after the shock, pointwise behaviour can be discontinuous, so $L^1$ estimates become more informative than pointwise bounds. A nonlinear interaction can remain small in an integral norm even when its pointwise behaviour changes dramatically at shock formation.
\end{explained}

